
\begin{figure*}[ht]
\centering
% ===== Row: Car 1 =====
\begin{tabular}{lcc}
  \centering
  % left: stacked range images
\raisebox{5mm}{\multirow{1}{*}{\rotatebox[origin=c]{90}{\footnotesize GT}}} & 
\includegraphics[width=0.46\linewidth]{figures/car1/r3dpa_inpainting_gt_1_cropped_v3.pdf} & 
\includegraphics[width=0.46\linewidth]{figures/car2/r3dpa_inpainting_gt_2_cropped.pdf}\\ 
  
% \small \shortstack{Gen \\ RI} & 
\raisebox{5mm}{\multirow{1}{*}{\rotatebox[origin=c]{90}{ \footnotesize Gen}}} & 
\includegraphics[width=0.46\linewidth]{figures/car1/r3dpa_inpainting_gen_1_cropped_v3.pdf} & 
\includegraphics[width=0.46\linewidth]{figures/car2/r3dpa_inpainting_gen_2_cropped.pdf} \\
    
& \cropcenter{0.22\linewidth}{\detokenize{figures/car1/gt.png}}  \cropcenter{0.22\linewidth}{\detokenize{figures/car1/gen.png}} & 
\cropcenter{0.22\linewidth}{\detokenize{figures/car2/gt.png}}
\cropcenter{0.22\linewidth}{\detokenize{figures/car2/gen.png}} \\
& \footnotesize GT PC {~~~~~~~~~~~~~~~~~~~~~~~~~~~~~} \footnotesize Gen PC & \footnotesize GT PC  {~~~~~~~~~~~~~~~~~~~~~~~~~~~~~}  \footnotesize Gen PC   \\


\end{tabular}
\caption{\textbf{Test-time conditioning with 3D representations for object inpainting.} A user-specified mask (red) is drawn on top of the ground truth range image (center-cropped for visibility). The real scene is transported to the Gaussian distribution, then the user defined area is filled with noise. During the denoising process only the masked area is supervised with features corresponding to the target object class. The scenes are also presented projected to 3D points, colored according to the intersection with the frustum of the projected mask.}
\label{fig:inpaint}
\vspace{-0.2cm}
\end{figure*}